\scanpage{11}
\inserted{fonction \uline{réelle} mesurable, sur \(M\) \deleted{\unsure{telle}} \((|M|<+\infty)\)} telle \deleted{que \(\mu_f\) \unsure{étant quelconque positive il suffit que} \(\mu_f\) soit \unsure{bornée, donc que par exemple} \(|f^{-1}(\complement\{0\})|<+\infty\). Alors} il existe donc une fonction décroissante \(\varphi_f\) et une seule \inserted{sur \(]0,|M|[\)} qui soit spectralement équivalente à \(f\). On l'appelle encore le \uline{réarrangement décroissant} de la fonction réelle \(f\). Le résultat suivant est classique.

\uline{Proposition 2} Soient \(f,g\) deux fonctions réelles mesurables sur \(M\), \deleted{telles}\inserted{\((|M|<+\infty)\)} \(\bigl|\) Pour que \(fg\) soit sommable, il faut et il suffit que \(\varphi_f\varphi_g\) le soit, et on a alors
\[
(17)\qquad\int fg\,dm\leq\int_0^{|M|}\varphi_f\varphi_g\,ds
\]

\deleted{\unsure{Ex.}}

\mspage{12}
